problem:  
Let $\{(M_i,g_i)\}$ be a sequence of closed, non‑collapsing $n$‑manifolds with $|{\rm Sec}_{g_i}|\le1$ and $\operatorname{Vol}(M_i)\ge v_0>0$, $\operatorname{diam}(M_i)\le D$, converging to $(X,d_X)$ in the Gromov–Hausdorff sense.  
Let $G_2(TM_i)$ denote the Grassmann bundle of \(2\)-planes with normalized measure $\nu_i$, and let \(K_i(p,\sigma)={\rm Sec}_{g_i}(\sigma)\).  Define the induced probability measure
\[
\mu_i=(K_i)_*\nu_i\in\mathcal P([-1,1]).
\]
We already know that a subsequence of $\mu_i$ converges weakly.  

**Refined Problem (Set of Curvature Distribution Limits):**  
Describe the **set of all possible weak limits**
\[
\mathcal L(X)=\{\;\mu_\infty:\exists (M_i,g_i)\!\to\! X,\ |{\rm Sec}_{g_i}|\le1,\ \operatorname{Vol}(M_i)\!\ge\!v_0\;\}
\subset\mathcal P([-1,1]).
\]
1. Does $\mathcal L(X)$ depend only on $(X,d_X)$?  
2. For which limit spaces \(X\) (e.g. smooth \(C^{1,\alpha}\) manifolds, Ricci‑limit spaces, or convex Alexandrov spaces) is $\mathcal L(X)$ a **singleton**?  
3. In general, what geometric features of $X$ are detectable from $\mathcal L(X)$—for instance, do extremal points of $\mathcal L(X)$ correspond to different smoothing regimes of $X$?  

Why is it a "good" problem:  
This formulation avoids the unrealistic uniqueness conjecture yet still asks a deep and well‑posed question about the *space of possible curvature distributions* determined by the limit geometry. It investigates how fine curvature information degenerates under Gromov–Hausdorff convergence and seeks to classify to what extent the limit space constrains higher‑order geometric data of its smooth approximations.  

Such a problem lies on the frontier between differential and synthetic geometry: it would clarify what remains of directional curvature information once one passes to a purely metric notion of limit. A positive structural description of $\mathcal L(X)$ would create a novel invariant for Alexandrov or Ricci‑limit spaces and might connect to curvature measures and metric stability theory—offering a compact, elementary statement that conceals profound geometric depth.